\documentclass[UTF8,aspectratio=169]{ctexbeamer}
\usetheme{Madrid}
\usecolortheme{default}
\usepackage{amsmath}
\usepackage{booktabs}
\usepackage{siunitx}

\title{Toy Inference From a Saved Run Directory}
\author{Mie Spray Postprocessing Project}
\date{\today}

\begin{document}

\begin{frame}
  \titlepage
\end{frame}

\begin{frame}{Audience and Goal}
\textbf{Audience:} Engineers who want a quick post-training sanity check without rebuilding the full dataset.

\vspace{0.5em}
\textbf{Goal of this talk:}
\begin{itemize}
  \item Explain what \texttt{toy\_inference\_from\_run.ipynb} does.
  \item Show how manual physical inputs are transformed into model-compatible features.
  \item Interpret the notebook outputs: mean penetration, uncertainty band, and simple collision proxies.
  \item Clarify what is useful now and what is still only a geometric approximation.
\end{itemize}
\end{frame}

\section{Part 1: Inference Workflow}

\begin{frame}{1.1 Notebook Workflow}
\textbf{Core workflow:}
\begin{enumerate}
  \item Set \texttt{RUN\_DIR} and manual physical inputs.
  \item Load \texttt{train\_config\_used.json}, \texttt{scaler\_state.json}, and model checkpoint.
  \item Rebuild the exact training-time feature scaling.
  \item Sweep time from \SI{0}{ms} to \SI{5}{ms} and predict $\mu(t)$ and $\log \sigma^2(t)$.
  \item Visualize uncertainty and project the result into a simple 2D chamber geometry.
\end{enumerate}
\end{frame}

\begin{frame}{1.2 Run Artifacts Used by the Notebook}
\textbf{Observed default run directory in the notebook:}
\begin{itemize}
  \item \texttt{MLP/runs\_mlp/stage2\_NLL\_penetration\_20260317\_194155}
\end{itemize}

\vspace{0.3em}
\textbf{Artifacts loaded from that run:}
\begin{itemize}
  \item network architecture and hyperparameters from \texttt{train\_config\_used.json}
  \item scaling statistics from \texttt{scaler\_state.json}
  \item model weights from the best available checkpoint
\end{itemize}

\vspace{0.3em}
\textbf{Observed notebook output:}
\begin{itemize}
  \item device: \textbf{cuda}
  \item feature dimension: \textbf{9}
  \item checkpoint found in the example run: \texttt{best\_model\_stage1.pt}
\end{itemize}
\end{frame}

\section{Part 2: Manual Inputs to Network Features}

\begin{frame}{2.1 Example Toy Input Condition}
\textbf{Default manual inputs in the notebook:}
\begin{itemize}
  \item tilt angle = $20^\circ$
  \item plumes = 10
  \item nozzle diameter = \SI{0.355}{mm}
  \item injection duration = \SI{1000}{\micro\second}
  \item injection pressure = \SI{2000}{bar}
  \item chamber pressure = \SI{25}{bar}
  \item control backpressure = \SI{4}{bar}
\end{itemize}

\vspace{0.3em}
\textbf{Time sweep:}
\begin{itemize}
  \item \num{300} points over \SIrange{0}{5}{ms}
\end{itemize}
\end{frame}

\begin{frame}{2.2 How Raw Inputs Are Scaled}
\textbf{The notebook reuses the same preprocessing logic as training:}
\begin{itemize}
  \item time is normalized to $[0,1]$ using the saved min/max
  \item pressure terms use $\log P_{\text{inj}}$, $\log P_{\text{ch}}$, and $\log(P_{\text{inj}}-P_{\text{ch}})$
  \item geometry and condition terms are z-scored using the saved train statistics
\end{itemize}

\vspace{0.3em}
\textbf{Resulting feature vector per time point:}
\[
x_t =
[
\text{time\_norm},
\text{tilt}_z,
\text{plumes}_z,
\text{diameter}_z,
\text{duration}_z,
\log P_{\text{inj},z},
\log P_{\text{ch},z},
\log \Delta P_z,
\text{backpressure}_z
]
\]

\vspace{0.3em}
This guarantees inference uses the same input space as training.
\end{frame}

\section{Part 3: Penetration and Uncertainty Output}

\begin{frame}{3.1 What the MLP Predicts}
The loaded model outputs two channels:
\[
[\mu(t), \log \sigma^2(t)]
\]

\vspace{0.3em}
\textbf{Post-processing in the notebook:}
\[
\sigma(t) = \max\!\left(\sqrt{e^{\log \sigma^2(t)}},\ \sigma_{\min}\right)
\]

\vspace{0.3em}
\textbf{Then the notebook plots:}
\begin{itemize}
  \item predicted mean penetration $\mu(t)$
  \item predictive interval $\mu(t)\pm \sigma(t)$
  \item an optional \SI{90}{mm} reference line
  \item where $\sigma_{\min}$ is the saved floor from \texttt{std\_clamp\_min}
\end{itemize}
\end{frame}

\begin{frame}{3.2 Example Numerical Output}
\textbf{Observed time-sweep output from the notebook:}
\begin{itemize}
  \item feature matrix shape: \textbf{(300, 9)}
  \item predicted std range: \textbf{6.27 mm to 15.15 mm}
\end{itemize}

\vspace{0.3em}
\textbf{Example single-time query from the same run at \SI{0.5}{ms}:}
\begin{itemize}
  \item predicted mean penetration $\mu \approx \SI{32.74}{mm}$
  \item predicted axial std $\sigma \approx \SI{10.18}{mm}$
\end{itemize}

\vspace{0.3em}
This is the simplest way to check whether the Stage-2 model produces plausible uncertainty magnitudes under manually chosen conditions.
\end{frame}

\section{Part 4: 2D Geometry Proxy}

\begin{frame}{4.1 Simple Mask-Based Overlap Solver}
After 1D inference, the notebook builds a coarse 2D geometry:
\begin{itemize}
  \item cylinder radius = \SI{120}{mm}
  \item cylinder height = \SI{60}{mm}
  \item cylinder head offset = \SI{50}{mm}
  \item grid size = \SI{1}{mm}
  \item wall mask starts at $x \ge \SI{120}{mm}$
\end{itemize}

\vspace{0.3em}
\textbf{A simple deterministic proxy is formed by:}
\begin{itemize}
  \item projecting the predicted penetration along the spray tilt direction
  \item turning the predicted mean length and $+1\sigma$ length into binary masks
  \item measuring overlap area with wall and cylinder masks
\end{itemize}
\end{frame}

\begin{frame}{4.2 Gaussian Collision Proxy}
The notebook then upgrades the simple mask idea to a rotated 2D Gaussian:
\[
p(u,v) \propto
\exp\!\left(
-\frac{1}{2}\left[
\left(\frac{u}{\sigma_{\text{axis}}}\right)^2 +
\left(\frac{v}{\sigma_{\perp}}\right)^2
\right]
\right)
\]

\vspace{0.3em}
\textbf{How the widths are chosen:}
\begin{itemize}
  \item axial spread: $\sigma_{\text{axis}} = \sigma$ from model output
  \item orthogonal spread: plume width $\approx \mu \tan(\theta/2)$
  \item notebook assumption: $\sigma_{\perp} \approx \text{width}/3$
\end{itemize}

\vspace{0.3em}
The Gaussian is rotated by the spray tilt angle and integrated over wall/cylinder masks to get probability mass.
\end{frame}

\begin{frame}{4.3 Example Geometry Result}
\textbf{Observed notebook outputs for the default test condition:}
\begin{itemize}
  \item simple wall collision area: \textbf{0}
  \item simple cylinder collision area: \textbf{0}
  \item $+1\sigma$ simple overlaps: also \textbf{0}
\end{itemize}

\vspace{0.3em}
\textbf{Gaussian-mask probability masses:}
\begin{itemize}
  \item wall mass $\approx 1.07 \times 10^{-20}$
  \item cylinder mass $\approx 6.09 \times 10^{-23}$
  \item any collision mass $\approx 1.08 \times 10^{-20}$
\end{itemize}

\vspace{0.3em}
For this operating condition, the toy case is effectively far from collision.
\end{frame}

\begin{frame}{4.4 What This Notebook Is Good For}
\textbf{Current strengths:}
\begin{itemize}
  \item fast sanity check for trained runs
  \item manual sensitivity study over pressure, duration, tilt, and diameter
  \item quick interpretation of predicted uncertainty in physical units
  \item simple geometry screening before building a more complete solver
\end{itemize}

\vspace{0.3em}
\textbf{Current limitations:}
\begin{itemize}
  \item 2D spread is assumed, not learned directly from images
  \item orthogonal Gaussian width is a heuristic from cone angle
  \item collision metric is only a proxy, not a CFD or multiphase transport model
\end{itemize}
\end{frame}

\begin{frame}{Takeaway}
\begin{itemize}
  \item \texttt{toy\_inference\_from\_run.ipynb} turns a saved run directory into a practical engineering demo tool.
  \item It exactly reuses the saved scaler and model interface, so toy inputs stay training-compatible.
  \item The notebook does more than draw $\mu \pm \sigma$: it also maps 1D uncertainty into a simple 2D collision proxy.
  \item This makes it a useful bridge between learned penetration prediction and geometry-aware downstream reasoning.
\end{itemize}
\end{frame}

\end{document}
